Η παρόν εργασία μελετά τον σχεδιασμό ενός ρομποτικού διασώστη. Ο διασώστης αυτός προορίζεται να χρησιμοποιηθεί σε αποστολές έρευνας και διάσωσης σε αστικές περιοχές (U.S.A.R.), όπου η πρόσβαση για ανθρώπους διασώστες μπορεί να είναι δύσκολη έως και πολύ επικίνδυνη. Ο ρομποτικός διασώστης θα είναι ικανός να κινείται σε διάφορα είδη εδάφους, να αποφεύγει εμπόδια, να εντοπίζει θύματα και να μεταφέρει βασικό εξοπλισμό πρώτων βοηθειών. Η μελέτη περιλαμβάνει τον προκαταρτικό και διαμορφωτικό σχεδιασμό του ρομπότ, καθώς και την ανάλυση των απαιτήσεων και προδιαγραφών του συστήματος. Στόχος της εργασίας είναι να παρουσιάσει μια ολοκληρωμένη προσέγγιση για τον σχεδιασμό ενός ρομποτικού διασώστη, ικανού να λειτουργεί ημιαυτόματα, λαμβάνοντας υπόψη τις προκλήσεις και τις απαιτήσεις που σχετίζονται με την αποστολή έρευνας και διάσωσης σε αστικές περιοχές. 
